\subsection{Task, Corpus, and Benchmark}\label{sec:benoit-task}

The document-label experiments instantiate the abstract oracle
$\fstar$ of Section~\ref{sec:theorems} as the rubric social scientists
actually use. The summarizer $g$ is a prompted language model, the
merge is an LLM-friendly prompt, and the readout $f$ is a second LLM
call under a scoring rubric that maps the root summary to a
seven-point scale. The target benchmark is
\citet{BenoitEtAl2025}, which establishes the current empirical bar
for LLM-based political-text measurement. Both $g$ and $f$ are fixed
open-weight checkpoints driven by rubric prompts: the optimization we
perform is over the prompts, chunk sizes, and audit design, never over
model weights.

The change of parameterization leaves the framework's objects intact.
The training objective is still $J_\lambda$ from
Equation~\eqref{eq:oracle-projection-objective}, minimized over prompts
and chunking rather than operator weights; the audit of
Section~\ref{sec:estimation} needs only realized leaf, merge, and
on-range calls plus logged sampling propensities, and is indifferent to
whether the parameters behind those calls are FNO weights or a decoding
temperature; and the gap bound is stated on $f$ and $g$, not on a
particular parameterization. The load-bearing deployed correctness step
is therefore the same audited operator interface as in the controlled
settings, now instantiated with an LLM summarizer and an LLM-scored
readout.

The long-document motivation is practical rather than cosmetic.
Million-token context windows reduce hard truncation, but they do not
by themselves certify that a single monolithic read preserved every
rubric-relevant distinction, nor do they give the researcher a local
failure mode when it did not. The tree makes that claim auditable: if
sufficiently small leaves and local merges preserve the policy
evidence, then the root score should be stable as the chunk schedule
changes, and the experiments below use that stability as the
diagnostic. We ask whether a single open-weight model driven through
the same local laws can match or exceed Benoit's $18$-score ensemble
of three proprietary frontier LLMs on the same Manifesto Project
benchmark. The short answer is yes: at an $8\mathrm{K}$-character
leaf, C-TreePO on Gemma-4-31B-NVFP4 produces a macro Pearson $r$ of
$0.829$ against expert-survey means, vs.\ $0.817$ for Benoit's
proprietary ensemble and $0.793$ for the matched Gemma-3-27B
open-weight replication, at roughly two orders of magnitude lower
compute.

\paragraph{The rubric oracle.} The oracle is the set of scoring
rubrics standardized by \citet{BenoitEtAl2025}, which place
every party platform on a seven-point scale for each of six policy
dimensions: economic (public services vs.\ taxation), social
(liberal vs.\ conservative), immigration, European integration,
environment, and decentralization. For a span $x$ of manifesto text and
dimension $d$, $\fstar_d(x) \in \{1,\dots,7\}$ is the rubric-indexed
judgment a trained coder would assign; the full oracle is the $6$-vector
$(\fstar_1(x),\dots,\fstar_6(x))$. This is rubric-driven semantics, not
a closed-form algorithm. Its computation requires a human-readable
scoring context and either an expert coder or a language model
instructed by that context. There is no hand-designed mergeable sketch
for it, which is precisely the regime the framework exists to cover.

\paragraph{Data and benchmark.} The comparison unit is a
manifesto--dimension observation. A manifesto
is a party's election platform from the Manifesto Project corpus; a
dimension is one of the six seven-point policy scales above. Benoit et
al.'s main manifesto benchmark contains $235$ party platforms from
parliamentary democracies, matched to election years for which expert
survey estimates of party position are available. Their broader paper
also reports $23$ Kl\"uver coalition agreements, but the headline
comparison in Table~\ref{tab:benoit-headline} is organized around the
Manifesto Project platform benchmark.

The target value for each manifesto--dimension pair is an expert-survey
mean, not a single coder's label. Benoit et al.'s released benchmark
provides one harmonized mean party-position estimate on the seven-point
scale for each platform and policy dimension. We evaluate a method by
correlating its predicted scores with those expert means across the
platforms in a dimension; the macro score is the unweighted mean of the
six dimension-level Pearson correlations.

One additional row in Benoit et al.'s paper matters for interpreting
these correlations. Their Table~3 reports a split-expert agreement
reference: within each country in the Benoit--Laver replication survey,
experts are randomly divided into two groups, each group's mean
party-position estimate is computed on the same issue-position scale,
and the two resulting vectors of mean estimates are correlated. This is
an empirical agreement reference for expert-survey measurement, not a
theoretical limit and not a statistic estimated on our held-out
split. It ranges from $0.78$ to $0.95$ Pearson $r$ across the six
dimensions, with macro $0.873$; the economic value used later in the
ladder figure is $r = 0.880$.

Benoit et al.'s LLM protocol gives the method we compare against. Each
manifesto--dimension pair runs through a two-stage pipeline: a summarizer
produces a $300$--$400$ word English summary, then a scorer reads that
summary and emits an integer on the seven-point scale. Their headline
ensemble uses
$3 \text{ LLMs} \times 2 \text{ shot-settings} \times 3 \text{ summaries}
= 18$ scores per manifesto--dimension. Three published numbers therefore
matter for what follows: Table~3 gives the split-expert reference just
defined; Figure~1 gives the proprietary ensemble headline ($0.49$--$0.92$
across dimensions, macro $0.817$); and Table~6 gives the open-weight
Gemma-3-27B-IT replication (macro $0.793$). A contender must reproduce
the per-dimension correlations, not just the macro.

\paragraph{Plugging the C-Tree into the Benoit corpus.} We reuse
Benoit's expert-mean benchmarks and manifesto crosswalks
without modification, so the results below are directly stackable next
to the numbers just cited. The only substitution is the pipeline. A
manifesto is partitioned into contiguous character spans of size
$c\in\{4\mathrm{K},8\mathrm{K},16\mathrm{K},24\mathrm{K},32\mathrm{K},64\mathrm{K}\}$;
each span is summarized by a single shared language model
(\textsc{Gemma-4-31B-IT-NVFP4}) under a rubric that names the
preservation categories drawn from Benoit's scoring contexts; adjacent
summaries are merged pairwise under the same rubric up to a root;
and the root summary is then scored on the seven-point scale. Two
pipeline variants are reported:

\begin{itemize}[leftmargin=1.5em]
    \item \textbf{Per-dimension C-Tree.} One summarizer and one scorer
          per dimension; six independent trees per manifesto. The
          rubric at every node names only that dimension's preservation
          categories. This is the setting the theorems of
          Section~\ref{sec:theorems} were stated for.
    \item \textbf{Joint C-Tree.} A single summarizer per manifesto
          with a union rubric covering all six dimensions; the
          resulting root summary is scored six times, once per
          dimension. The scoring stage is reordered so all $N$
          scoring calls for one dimension execute consecutively,
          which lets the per-dimension system prompt be reused across
          calls and amortizes the dominant prefill cost.
\end{itemize}

Both variants close under the same three local laws as the controlled
settings: C1 (leaf sufficiency) requires each leaf summary to
preserve the rubric-relevant evidence in its span; C2 (idempotence)
requires re-summarizing a leaf summary to leave the predicted score
unchanged; C3 (Triangle / Merge Consistency) requires the ordered merge
summary to preserve the same rubric evidence as the corresponding raw
union span.
Order-swap or permutation probes are useful stress tests when the chosen
rubric treats quasi-sentence evidence as bag-like, but they are not the
core C3 law. In the language of Section~\ref{sec:theorems}, the prompt-only
pipeline is a semantic instantiation of the same operator-plus-projection
interface: oracle fit comes from the root score, while local-law weights
and audits target the theorem-eligible subspace. The joint variant adds
one invariant not stated as a fourth law: the same leaf summary must
support all six scorers. Empirically this invariant matters more than any
within-law tightening because it forces $g$ to compress without discarding
dimension-specific evidence.

\subsection{Document Labels Only: Headline Parity and Chunk Invariance}\label{sec:benoit-headline}

Table~\ref{tab:benoit-headline} reports per-dimension Pearson $r$
against Benoit's expert-survey means for the proprietary ensemble, the
matched Gemma-3-27B open-weight replication, and five C-TreePO
configurations on \textsc{Gemma-4-31B-IT-NVFP4}. Two rows carry the
claim.

\begin{itemize}[leftmargin=1.5em]
    \item \textbf{Per-dimension C-Tree, $c = 8\mathrm{K}$ leaves.}
          Macro $0.829$, above the proprietary ensemble's $0.817$ and
          $+0.036$ above the matched open-weight replication. The
          three dimensions with the biggest gain against Gemma-3
          are decentralization ($+0.093$), economic ($+0.079$), and
          EU ($+0.061$). The losses are environment ($-0.021$) and
          immigration ($-0.006$).
    \item \textbf{Per-dimension C-Tree, $c = 4\mathrm{K}$ leaves}
          (tiny-leaf extension). Macro $0.842$, within $0.031$ of
          Benoit's \emph{split-expert} reference ($0.873$). The
          two dimensions where we exceed that empirical reference are
          EU and decentralization, both flagged by Benoit as
          low-variance boundary cases in their Figures~2--3.
\end{itemize}

\begin{table}[htbp]
    \centering
    \resizebox{\linewidth}{!}{% Generated table; do not edit by hand.
% Metric: Pearson r (higher better)
\begin{tabular}{lrrrrrrrr}
\toprule
Method & Econ. & Social & Immig. & EU & Env. & Decent. & Macro & \#/6 \\
\midrule
\multicolumn{9}{l}{\textit{Benoit et al. (2025) reference}} \\
Proprietary ensemble, 18 scores (Fig 1) & $+0.870$ & $+0.920$ & $+0.890$ & $+0.910$ & $+0.820$ & $+0.490$ & $+0.817$ & 6/6 \\
Split-expert reference (Table 3) & $+0.880$ & $+0.910$ & $+0.880$ & $+0.950$ & $+0.840$ & $+0.780$ & $+0.873$ & 6/6 \\
LLaMA-3.3-70B (Table 6) & $+0.840$ & $+0.870$ & $+0.860$ & $+0.860$ & $+0.680$ & $+0.400$ & $+0.752$ & 6/6 \\
DeepSeek-V3 (Table 6) & $+0.840$ & $+0.870$ & $+0.890$ & $+0.860$ & $+0.790$ & $+0.450$ & $+0.783$ & 6/6 \\
Gemma-3-27B-IT (Table 6) & $+0.860$ & $+0.860$ & $+0.890$ & $+0.840$ & $+0.860$ & $+0.450$ & $+0.793$ & 6/6 \\
\midrule
\multicolumn{9}{l}{\textit{Ours: per-dim pipeline $\times$ leaf size (Gemma-4-31B-NVFP4, 1 summarizer + 1 scorer per dim)}} \\
leaf = 64 K chars ($\approx$16K tokens) & $+0.916$ & $+0.817$ & $+0.888$ & $+0.924$ & $+0.854$ & $+0.489$ & $+0.815$ & 6/6 \\
leaf = 32 K chars ($\approx$8K tokens) & $+0.929$ & $+0.857$ & $+0.886$ & $+0.916$ & $+0.877$ & $+0.480$ & $+0.824$ & 6/6 \\
leaf = 24 K chars ($\approx$6K tokens) --- full test n$\approx$215 & $+0.892$ & $+0.840$ & $+0.867$ & $+0.896$ & $+0.814$ & $+0.464$ & $+0.796$ & 6/6 \\
leaf = 16 K chars ($\approx$4K tokens) & $+0.925$ & $+0.847$ & $+0.898$ & $+0.909$ & $+0.831$ & $+0.537$ & $+0.824$ & 6/6 \\
leaf =  8 K chars ($\approx$2K tokens) & $+0.939$ & $+0.869$ & $+0.884$ & $+0.901$ & $+0.839$ & $+0.543$ & $+0.829$ & 6/6 \\
\midrule
\multicolumn{9}{l}{\textit{Ours: combined pipeline $\times$ leaf size (one shared summarizer with union rubric, 6 scores)}} \\
leaf = 64 K chars & $+0.910$ & $+0.867$ & $+0.897$ & $+0.910$ & $+0.795$ & $+0.458$ & $+0.806$ & 6/6 \\
leaf = 32 K chars & $+0.917$ & $+0.849$ & $+0.908$ & $+0.889$ & $+0.749$ & $+0.438$ & $+0.792$ & 6/6 \\
leaf = 24 K chars (full test n=229) & $+0.868$ & $+0.850$ & $+0.880$ & $+0.902$ & $+0.809$ & $+0.396$ & $+0.784$ & 6/6 \\
leaf = 16 K chars & $+0.919$ & $+0.851$ & $+0.916$ & $+0.889$ & $+0.808$ & $+0.441$ & $+0.804$ & 6/6 \\
leaf =  8 K chars & $+0.927$ & $+0.874$ & $+0.911$ & $+0.917$ & $+0.801$ & $+0.413$ & $+0.807$ & 6/6 \\
\midrule
\multicolumn{9}{l}{\textit{Ours: tiny-leaf extensions of the chunk sweep (Gemma-4)}} \\
tree, leaf = 4 K chars & $+0.932$ & $+0.886$ & $+0.885$ & $+0.931$ & $+0.838$ & $+0.580$ & $+0.842$ & 6/6 \\
tree, leaf = 2 K chars & $+0.924$ & $+0.867$ & $+0.887$ & $+0.906$ & $+0.793$ & $+0.584$ & $+0.827$ & 6/6 \\
tree, leaf = 1 K chars (stress test, depth $\geq$6) & $+0.933$ & $+0.859$ & $+0.914$ & $+0.925$ & $+0.838$ & $+0.493$ & $+0.827$ & 6/6 \\
\midrule
\multicolumn{9}{l}{\textit{Ours: concat-no-merge (chunks summarized independently, joined, scored --- tests whether the merge step carries signal)}} \\
concat, leaf = 32 K chars & $+0.929$ & $+0.856$ & $+0.897$ & $+0.920$ & $+0.848$ & $+0.571$ & $+0.837$ & 6/6 \\
concat, leaf = 16 K chars (default) & $+0.931$ & $+0.823$ & $+0.909$ & $+0.912$ & $+0.866$ & $+0.523$ & $+0.827$ & 6/6 \\
concat, leaf =  8 K chars & $+0.932$ & $+0.849$ & $+0.882$ & $+0.915$ & $+0.850$ & $+0.594$ & $+0.837$ & 6/6 \\
\midrule
\multicolumn{9}{l}{\textit{Ours: flat baseline (no chunk, no summary; truncate text and score) --- tests whether tree is needed at all}} \\
flat, truncation = 48 K chars & $+0.936$ & $+0.821$ & $+0.926$ & $+0.889$ & $+0.806$ & $+0.587$ & $+0.827$ & 6/6 \\
flat, truncation = 24 K chars (default) & $+0.929$ & $+0.847$ & $+0.889$ & $+0.882$ & $+0.829$ & $+0.542$ & $+0.820$ & 6/6 \\
flat, truncation = 12 K chars & $+0.929$ & $+0.844$ & $+0.938$ & $+0.880$ & $+0.779$ & $+0.564$ & $+0.822$ & 6/6 \\
flat, truncation =  6 K chars & $+0.926$ & $+0.811$ & $+0.967$ & $+0.881$ & $+0.886$ & $+0.243$ & $+0.786$ & 6/6 \\
\midrule
\multicolumn{9}{l}{\textit{Ours: scorer ablations (Gemma-4, Benoit's GPT-4o summaries held fixed)}} \\
1. Zero-shot scorer & $+0.832$ & $+0.886$ & $+0.858$ & $+0.905$ & $+0.668$ & $+0.311$ & $+0.743$ & 6/6 \\
2. Few-shot optimized scorer & $+0.822$ & $+0.892$ & $+0.853$ & $+0.912$ & $+0.627$ & $+0.297$ & $+0.734$ & 6/6 \\
3. Joint scorer baseline (shared across 6 dims) & $+0.837$ & $+0.881$ & $+0.852$ & $+0.904$ & $+0.658$ & $+0.314$ & $+0.741$ & 6/6 \\
4. Joint scorer, few-shot optimized & $+0.817$ & $+0.891$ & $+0.848$ & $+0.908$ & $+0.700$ & $+0.313$ & $+0.746$ & 6/6 \\
5. Joint scorer, reflective prompt-optimized & $+0.832$ & $+0.882$ & $+0.848$ & $+0.879$ & $+0.688$ & $+0.344$ & $+0.746$ & 6/6 \\
\midrule
\multicolumn{9}{l}{\textit{Ours: exact-model replication (Benoit's Gemma-3-27B-IT BF16)}} \\
Gemma-3-27B scorer-only, OUR scoring rubric & $+0.826$ & $+0.864$ & $+0.846$ & $+0.902$ & $+0.573$ & $+0.329$ & $+0.723$ & 6/6 \\
Gemma-3-27B scorer-only, exact Benoit rubric & $+0.814$ & $+0.895$ & $+0.852$ & $+0.852$ & $+0.616$ & $+0.359$ & $+0.731$ & 6/6 \\
Gemma-3-27B, raw-prompt Benoit format (bare-integer response) & $+0.782$ & $+0.901$ & $+0.854$ & $+0.874$ & $+0.590$ & $+0.368$ & $+0.728$ & 6/6 \\
\bottomrule
\end{tabular}
}
    \caption{Pearson $r$ against Benoit expert-survey means across six
    policy dimensions. The first block reproduces
    \citet{BenoitEtAl2025}'s published numbers (Fig.~1 proprietary
    ensemble, Tab.~3 split-expert reference, Tab.~6 open-weight
    replication at Gemma-3-27B). The Table~3 row is an empirical
    agreement reference, not a theoretical limit. The remaining
    blocks are C-TreePO on Gemma-4-31B-NVFP4 across chunk sizes $c$,
    under the per-dimension
    and joint variants. The macro column is the unweighted mean
    across dimensions for each row. Full per-cell table with
    coverage, ablations, and compute accounting appears in
    Appendix~\ref{app:benoit-replication}.}\label{tab:benoit-headline}
\end{table}

Three structural statements follow from the table.

First, a single open-weight model matches a frontier ensemble. The per-dimension row at $c = 8\mathrm{K}$ (macro $0.829$) exceeds a three-frontier-LLM ensemble with $18$ scores per manifesto--dimension (macro $0.817$), on one open-weight model with one score per manifesto--dimension. The ensemble effect that \citet{BenoitEtAl2025} emphasize is therefore not necessary for open-weight replication on this corpus. A tree over a single model recovers the same macro-level agreement; the per-dimension dispersion tracks theirs to within sampling noise on five of six dimensions.

The flat and concat rows are controls for the same preservation claim.
The flat row shows that the expert-survey signal is visible to a direct
long-context read when enough manifesto text fits in context (macro
$0.827$ at a $48\mathrm{K}$-character truncation). The concat rows show
that the signal also survives a first distribution over leaves:
concatenated leaf summaries reach macro $0.837$ at both $32\mathrm{K}$
and $8\mathrm{K}$ leaves. The C-Tree rows complete the measurement
story. They keep the same signal after the document is broken into
smaller local units, with macro $0.829$ at $8\mathrm{K}$ and $0.842$ at
$4\mathrm{K}$, and give the audit many leaf and merge sites to sample.
That is the verification advantage: accuracy is preserved while the
evidence is spread over cheaper, more numerous, locally checkable
measurement points.

Second, leaf-size invariance is the main structural result. Across chunk
sizes varying by $16\times$ ($c = 4\mathrm{K}$ to $c = 64\mathrm{K}$),
macro $r$ varies by only $0.027$ in the per-dimension variant. The root
score is therefore not an artifact of one chunk schedule. In the theorem,
this is Proposition~\ref{prop:mergeable-reduction}'s prediction for a
semantic oracle: under local-law validity the root-level score does not
depend on how the document is partitioned. In the applied pipeline, it
means we can choose leaves small enough to make measurement easy to
sample, supervise, and audit without changing the estimand. Decentralization
is the instructive exception: its correlation \emph{improves} as leaves
shrink, from $0.489$ at $c = 64\mathrm{K}$ to $0.584$ at
$c = 2\mathrm{K}$. That is the dimension Benoit's own Figures~2--3 flag
as a boundary case, where coders and experts disagree about what local
evidence resolves a party's position. Finer localization helps exactly
where the relevant evidence would otherwise be compressed away. The
controlled counting benchmark shows the same signature
(Appendix~\ref{sec:markov-interlude}): when the oracle cares about
between-leaf interactions, more leaves help.

Third, the compute gap is roughly two orders of magnitude. The published \citet{BenoitEtAl2025} pipeline makes $18$ LLM calls per
manifesto--dimension, split across three proprietary frontier models,
for a corpus-wide total of $235 \times 6 \times 18 \approx 25{,}000$
frontier-API calls (plus an extra $23 \times 6 \times 18$ for the
Kl\"uver coalition agreements). The joint C-Tree variant makes
$2\lceil n_{\text{chars}}/c\rceil - 1 + 6$ calls per manifesto: one
summarize call per leaf, one merge call per internal node, and one
score call per dimension. For a typical manifesto with $c = 8\mathrm{K}$
leaves ($\approx 6$ leaves and $5$ internal merges) that is $17$ LM
calls against a single open-weight model on local hardware, versus
$18 \times 6 = 108$ frontier-API calls in Benoit's ensemble. At frontier
API list prices the compute gap is roughly two orders of magnitude per
manifesto; the full corpus runs reported in
Table~\ref{tab:benoit-headline} completed on a local GPU node inside a
working day. Appendix~\ref{app:benoit-replication} documents the
scoring-call reordering that Benoit's fixed-per-manifesto interleaving
prevents, and the wall-clock profile.

\subsection[Optimizing the Prompts: The Alternating f/g Ladder]{Optimizing the Prompts: The Alternating $f/g$ Ladder}\label{sec:rile-fg-ladder}

The headline table above freezes a single prompt for the summarizer/merge
program $g$ and a single prompt for the scorer/readout program $f$, then
sweeps only over chunk size. The iterative version is the
\emph{alternating $f/g$ prompt ladder}: optimize the scorer prompt against
the current summaries, then optimize the summarizer/merge prompt against
the current scorer, and repeat. This is the LLM-specific version of the
training story from the controlled settings. There, known oracle values
and local-law probes can supervise the operator and readout directly.
Here, the objects being changed are prompts, so co-adaptation itself has
to be measured: does repeated prompt tuning preserve chunk-invariance and
external agreement, or does one side simply overfit the current behavior
of the other?

Appendix~\ref{app:rile-fg-ladder} gives the full stage sequence and the
per-cell table. The compact label $f^a g^b$ means ``use the scorer after
$a$ scorer-side prompt optimizations and the summarizer after $b$
summarizer-side prompt optimizations''; it is a stage label, not
multiplication. That notation appears below only to identify the few
plotted cells discussed.

The ladder runs at two scopes: the economic dimension in depth under
two initialization lanes, and a single lane for each of the six
Benoit dimensions. The fully populated economic lane inherits
Benoit's archived GPT-4o summary as $g^0 = g^{\mathrm{Benoit}}$ and
starts from our scorer-only optimized prompt as $f^1$. Its plotted anchor is therefore
$f^1 g^0$, not a fresh $f^0 g^0$ cell. This is the natural continuation
of the fixed-summary scorer ablations in Appendix~\ref{app:benoit-replication},
and it lets the ladder ask what additional summarizer updates buy once
the scorer has already been tuned on Benoit's summary distribution. The
cleaner ``raw-init'' variant, where $(f^0, g^0)$ both come from the
own-Gemma baseline, is the setting we ultimately want the framework to
solve; both lanes are reported side-by-side in
Appendix~\ref{app:rile-fg-ladder}, and
Figure~\ref{fig:manifesto-fg-headline} below shows the main-text summary.

Two reference points frame the numbers below. The upper reference for
external Pearson is the economic value from \citet{BenoitEtAl2025}'s
Table~3, $r = 0.880$. As above, this is an agreement reference, not a
theoretical limit or a number estimated on our held-out split. The
lower reference for the
internal-external Pearson gap is $0$: the gap is internal Pearson
against the teacher trace minus external Pearson against expert means,
so values near zero say the scorer's internal agreement is not pulling
away from the external expert benchmark.

We report both lanes: the Benoit-init lane is fully populated across
the six leaf sizes, and the raw-init lane covers leaves
$256$--$8192$ tokens at the $f^1 g^0$ and $f^1 g^1$ stages (plus a few
deeper cells), enough to show the two lanes agree qualitatively on
the leaf-size-robustness claim.

The main empirical takeaway is leaf-size robustness: on the economic
dimension, the per-leaf best external Pearson sits at $r = 0.885$
(leaf 1024), $0.886$ (leaf 2048), $0.886$ (leaf 4096), $0.879$ (leaf
8192): a $0.007$-wide band ($[0.879, 0.886]$) straddling the
$r = 0.880$ split-expert reference across an $8\times$ range of
leaf sizes, with three of the four cells at or above it. This is the
divide-and-conquer claim in the LLM setting: fairly small leaves still
support the full-document measurement when the prompts are kept in
local summarizer/merge and scorer roles.
Below $\sim 1024$ tokens the method degrades gracefully
($r = 0.861$ at leaf $512$, $r = 0.830$ at leaf $256$), the
semantic analogue of the controlled benchmark's constraint that leaves
must be wide enough to carry the between-leaf evidence
(Appendix~\ref{sec:markov-interlude}). The same chunk-invariance pattern
Table~\ref{tab:benoit-headline} reports for the fixed-prompt pipeline
($\pm 0.027$ macro across $16\times$ leaf variation) is thus
reproduced here under iterative training. That is the non-obvious
outcome: medium-budget alternating prompt updates preserve the
fixed-prompt pipeline's chunk-insensitivity and leave the realized
operator near the split-expert reference rather than clearly below it.
Figure~\ref{fig:manifesto-fg-headline} plots the per-leaf best external
Pearson for both lanes against the split-expert reference.

\begin{figure}[t]
    \centering
    \includegraphics[width=0.72\linewidth]{manifesto_fg_ladder_headline.pdf}
    \caption{\textbf{Leaf-size robustness of the alternating $f/g$ ladder.}
    Per-leaf best external Pearson against Benoit expert means on the
    economic policy dimension, for the Benoit-init lane
    ($g^0 = g^{\mathrm{Benoit}}$, solid) and the raw-init lane ($g^0$ =
    own-Gemma baseline, dashed). The dash-dot reference at $r = 0.880$ is
    the \citet{BenoitEtAl2025} Table~3 economic split-expert correlation
    (an agreement reference, not a theoretical ceiling). Both lanes
    converge into a plateau at $r \approx 0.88$ for leaves
    $\geq 1024$ tokens; the raw-init lane additionally attains
    $r = 0.909$ at leaf $256$, stage $f^1 g^1$, with the tightest
    internal-external Pearson gap in the grid ($0.056$). See
    Appendix~\ref{app:rile-fg-ladder} for the full grids.}
    \label{fig:manifesto-fg-headline}
\end{figure}

The plateau costs nothing on the local-law probes the audit of
Section~\ref{sec:estimation} relies on: across every populated
Benoit-init cell, internal Pearson stays in $[0.953, 0.990]$ and the
internal-external Pearson gap stays in $[0.086, 0.161]$. The ladder's
best Benoit-init cell reaches $r = 0.886$ on economic, a $+0.049$
improvement over the strongest scorer-only ablation ($r = 0.837$) and
$+0.064$ over the few-shot scorer-only ablation ($r = 0.822$;
Appendix~\ref{app:benoit-replication}). Those baselines optimize only
$f$ while holding the summary distribution fixed, so the remaining gain
is the summarizer-update half of the loop. Comparison against the
per-dimension headline row of Table~\ref{tab:benoit-headline} is out
of scope: the ladder's leaf axis is in tokens, the headline's is in
characters, and no fixed-prompt sweep at matched token leaves and
prompt-search budget exists.

A raw-init variant where $(f^0, g^0)$ are both drawn from the
own-Gemma baseline pipeline (no Benoit summaries) tracks the
Benoit-init plateau closely for leaves $\geq 1024$ tokens
(Figure~\ref{fig:manifesto-fg-headline}, dashed series). Its one
departure from the Benoit-init shape is at the smallest leaf, where
the single fully-updated cell (leaf $256$, $f^1 g^1$) reaches
$r = 0.909$ (above the $0.880$ split-expert reference on this
held-out split), with the tightest internal-external Pearson gap
of the grid ($0.056$). Appendix~\ref{app:rile-fg-ladder} reports
the completed cells for both lanes.

The same ladder covers all six Benoit dimensions. One
single-dimension lane per dimension, over token leaves $256$--$8192$
and stages $f^1 g^0$ through $f^3 g^3$ ($36$ dimension--leaf cells,
all complete), reaches best held-out external Pearson $0.897$ on
economic (leaf $4096$), $0.928$ on social ($2048$), $0.894$ on
immigration ($4096$), $0.971$ on European integration ($1024$),
$0.779$ on environment ($8192$), and $0.540$ on decentralization
($256$). Five of six dimensions reach the neighborhood of the
$0.78$--$0.95$ band of Benoit's Table~3 split-expert references,
with environment at the band's lower edge;
decentralization stays the hard case at every leaf size and stage,
matching the chunk-sweep pattern of Table~\ref{tab:benoit-headline}.
The prompt updates help unevenly: economic, social, and European
integration peak two or three alternation steps past the anchor,
immigration and environment peak at the first step ($f^1 g^1$), and
decentralization peaks at the $f^1 g^0$ anchor itself, before any
summarizer update. Leaf-size robustness also recurs per dimension:
across the $32\times$ leaf range, the per-leaf best sits in a band
at most $0.056$ wide on economic, social, immigration, and European
integration. Appendix~\ref{app:rile-fg-ladder} reports the
per-dimension grid (Table~\ref{tab:fg-ladder-perdim}), the
leaf-sweep figures, and a joint variant with one shared summarizer,
which stays close to the single-dimension bests on five dimensions
and degrades decentralization.

\paragraph{Joint training and the sixth dimension.} The joint
variant's failure on decentralization is worth opening up, because it
is the clearest empirical picture in the corpus of what happens when
an optimization objective does not price the signal a dimension needs.
The joint $g$-step is rewarded by the unweighted mean of the six
per-dimension scorer rewards. At leaf $8096$ tokens the
decentralization trajectory oscillates: the first summarizer update
lifts external Pearson from $0.361$ to $0.461$, the largest
single-step gain anywhere in the joint grid; the next scorer update
pulls it back to $0.343$, below its own anchor; and the remaining
stages wander in the $0.34$--$0.41$ band without recovering the peak.
Environment pays for the excursions ($0.782$ at the anchor, $0.704$ by
the final stage), while the other four dimensions move by
$0.01$--$0.03$. The mechanism is multi-task interference through the
averaged reward. Decentralization is far from the consensus summary on
this corpus, so a summarizer update that serves it costs reward on the
other five dimensions, and the next scorer update, trained on the same
averaged metric, weights the decentralization signal at one sixth and
pulls the pair back toward the consensus. The alternation converts a
fixed disagreement between objectives into an oscillation rather than
a fixed point. The bottleneck is demonstrably on the scorer side: the
summarizer's output budget (twice the leaf size, enforced as a hard
context guard rather than silent truncation) is sufficient, since the
$f^1 g^1$ excursion shows $g$ can encode the decentralization evidence
within it.

The single-dimension lanes are the control on this diagnosis: trained
on decentralization alone, the pullback dynamic is absent, and the
dimension holds $0.540$ at leaf $256$ against the joint variant's
best of $0.461$. The practical recipe that falls out is two-track.
Default to the joint summarizer, which is cheaper (one tree per
manifesto) and matches the single-dimension bests within $0.044$ on
five of six dimensions; and watch the per-dimension $f^{*}$-gap, the
distance between internal agreement and external agreement, as the
fall-back trigger. In the joint run, the five well-served dimensions
keep that gap below $0.23$ at every stage (mostly below $0.16$), while
decentralization's never falls below $0.33$. A dimension whose gap sits
far outside the pack is being priced out of the shared objective, and
the fix is a single-target summarizer lane for that dimension, not
more joint iterations. Appendix~\ref{app:rile-fg-ladder} tabulates
the full trajectories behind these numbers.

\subsection{Preference Training on Oracle-Preserving Summaries}\label{sec:benoit-preference}

Theorem~\ref{thm:pref-equiv} makes the headline operationally useful
beyond correlation. When downstream preference training targets the
same oracle (``which summary better captures the manifesto's
position on dimension $d$?''), preferences can be collected on root
summaries $Z^{(R)}$ rather than on full documents $X$. Under local
preservation and the context-compatibility condition of
Assumption~\ref{ass:context}, the summary-based and
document-based DPO and GRPO population minimizers coincide; when
preservation is approximate, Theorem~\ref{thm:unified-gap} turns the
realized-violation estimate $\Delta_R$ into a quantitative drift
envelope on the training objective. The relevant side assumption is
oracle-indexed preferences:

\begin{assumption}[Oracle-indexed preferences]\label{ass:pref}
Given manifestos $x_1, x_2$ with oracle values
$y_1 = \fstar(x_1)$, $y_2 = \fstar(x_2)$, the probability that a human
prefers response $a_1$ to $a_2$ depends only on $(y_1, y_2, a_1, a_2)$,
not on the raw texts.
\end{assumption}

This is an application assumption package, not a theorem about manifesto
data. For DPO it instantiates oracle-measurable policy and
oracle-indexed pair-generation premises; for GRPO it instantiates the
analogous oracle-measurable policy/reward and oracle-indexed group
generation premises. The formal theorems then apply conditionally under
those premises. When preferences are noisy or only approximately
oracle-indexed, the residual is carried through the same
oracle-measurement term used in the quantitative certificates.

For a content-based rubric this is plausible: if two manifestos share
the same six-dimensional score and the same prompt, an ideally
calibrated annotator's preference depends on content, not surface
style. It can fail when annotators are swayed by tone, register, or
formatting, the usual DPO-data pathologies \citep{RafailovEtAl2023DPO}.
In that case Theorem~\ref{thm:unified-gap} replaces an equality with a
finite-sample envelope, and the operational recommendation is either
to widen the oracle to include the relevant stylistic features or to
report the certificate and let it carry the sensitivity.

The workflow implied by the theorems is to build hierarchies with a chosen chunk schedule (we settle on $c = 8\mathrm{K}$ as a default; the tiny-leaf extension at $c = 4\mathrm{K}$ is the cheapest path to the expert-reference line on the two boundary dimensions), audit with the sampling design of Section~\ref{sec:estimation} to bound $\Delta_R$, collect preferences on summaries rather than on full documents, train DPO or GRPO on the resulting (summary, preference) pairs, and deploy, reporting both the downstream estimate and the audit-based preference-gap certificate.
